\begin{resumen}
La transición hacia procesos de fabricación sostenibles ha impulsado el interés en la sinterización solar de arena, una técnica de impresión 3D que emplea luz solar concentrada para fundir sílice sin requerir energía eléctrica ni aglutinantes. Para viabilizar este proceso, resulta indispensable automatizar la orientación óptica, compensando el movimiento continuo del sol para mantener una estricta perpendicularidad y una distancia focal constante sobre el área de impresión.

El objetivo principal de este Trabajo de Fin de Grado es evaluar, diseñar y validar cinemáticamente arquitecturas robóticas de seguimiento solar en tiempo real. Se busca desarrollar un sistema gobernado por un bucle de control de lazo cerrado, utilizando visión artificial para garantizar la concentración térmica precisa mediante una lente de Fresnel solidaria al mecanismo.

La metodología aplicada ha comenzado con el desarrollo de un algoritmo modular de procesamiento de imagen que cuantifica el error de alineación óptica. Posteriormente, se ha evaluado una plataforma robótica paralela Gough-Stewart de seis grados de libertad. Mediante modelado matemático con matrices de rotación tridimensionales y cosimulación, se resolvió su cinemática inversa y se mapeó su espacio de trabajo angular. Estos resultados se cruzaron con las efemérides solares específicas de la latitud de Toledo, correlacionando la capacidad de la plataforma frente a los requerimientos de elevación y acimut en solsticios y equinoccios.

El análisis reveló limitaciones críticas: el espacio de trabajo angular de la plataforma Gough-Stewart queda restringido aproximadamente a ±20° en elevación y ±25° en acimut por singularidades mecánicas y colisiones entre actuadores, un rango insuficiente frente a los 240° acimutales y los más de 70° cenitales que el sol recorre durante el solsticio de verano en la latitud de Toledo. Esta incompatibilidad impide abarcar las trayectorias solares extremas exigidas a lo largo del año. Como respuesta, se sintetizó una segunda arquitectura desacoplada. 

Para la variación acimutal, se desarrolló un mecanismo plano de lazo cerrado de seis barras (variante de la cadena de Stephenson tipo III), accionado por un único actuador lineal. Este mecanismo alcanza un barrido de 240°, superando significativamente el rango angular solar máximo en este plano. Para la variación cenital, se sintetizó un segundo mecanismo independiente montado sobre el bastidor del primero. Se trata de una configuración biela-manivela en la que un actuador lineal impone la rotación de la estructura portadora de la lente de Fresnel alrededor de un eje horizontal solidario al subconjunto acimutal. Este mecanismo proporciona un barrido cenital efectivo de 80°, comprendido entre la posición horizontal (sol próximo al horizonte) y la vertical (cenit), abarcando el intervalo de horas solares de mayor irradiancia útil para el proceso de sinterización.

Finalmente se materializa un sistema de control para la arquitectura definitiva repartiendo las funciones entre tres nodos comunicados por una red WiFi local autónoma. Un ESP32-CAM actúa como punto de acceso y sensor del lazo: captura las imágenes solares y publica el vector de error del centroide mediante peticiones HTTP. Un segundo ESP32, conectado como cliente, asume la estimación y el control: acumula el error sobre el ángulo absoluto de la plataforma, resuelve la cinemática inversa por bisección numérica y comanda los actuadores por posición absoluta, evitando la acumulación de error de los esquemas incrementales. Una consola Python sobre PC opera como estación de telemetría sin intervenir en el lazo. Esta arquitectura permite al seguidor operar de manera totalmente autónoma, sin conexión a internet ni dependencia del PC.

En conclusión, este trabajo demuestra la inviabilidad cinemática de la plataforma Gough-Stewart clásica para el seguimiento solar ininterrumpido en la latitud estudiada. Sin embargo, propone y valida teóricamente una arquitectura híbrida optimizada que, combinando mecanismos de lazo cerrado y retroalimentación visual, logra amplificar el rango rotacional necesario para operar de forma eficiente y rígida durante el ciclo solar completo. La integración de este diseño mecánico con un bucle de control distribuido consolida un sistema autónomo, robusto y trasladable a otros escenarios donde la precisión de apuntamiento sea crítica.


\end{resumen}

\begin{abstract}
The push toward sustainable manufacturing has renewed interest in solar sand sintering, a 3D-printing technique that fuses silica using concentrated sunlight, without electrical power or chemical binders. Making this process practical hinges on automating the optical pointing, tracking the sun's continuous motion to keep the focal spot perpendicular and at a constant distance from the printing bed.

The main goal of this Bachelor's Thesis is to assess, design and kinematically validate robotic architectures for real-time solar tracking. The target system relies on a closed-loop control scheme driven by computer vision, ensuring accurate thermal concentration through a Fresnel lens rigidly mounted on the mechanism.

The work starts with a modular image-processing algorithm that quantifies the optical alignment error. A six-degree-of-freedom Gough-Stewart parallel platform is then evaluated: three-dimensional rotation matrices and a co-simulation environment are used to solve its inverse kinematics and map its angular workspace. These results are cross-checked against the solar ephemerides for the latitude of Toledo, matching the platform's reach against the elevation and azimuth demands at the solstices and equinoxes.

The analysis exposes a critical bottleneck: mechanical singularities and inter-leg collisions confine the Gough-Stewart workspace to roughly ±20° in elevation and ±25° in azimuth, a range that falls far short of the 240° of azimuth and the more than 70° of zenith the sun sweeps during the summer solstice in Toledo. The platform simply cannot reach the extreme solar trajectories that occur throughout the year. A second, decoupled architecture is therefore synthesised in response.

Azimuthal motion is handled by a planar six-bar closed-loop linkage (a Stephenson type III variant) driven by a single linear actuator. The mechanism delivers a 240° sweep, comfortably above the maximum solar azimuth range. Elevation is provided by an independent mechanism stacked on the azimuthal frame: a slider-crank arrangement in which a linear actuator rotates the Fresnel lens carrier about a horizontal axis fixed to the azimuthal subassembly. This second stage yields an effective zenithal sweep of 80°, spanning from horizon to zenith and covering the daylight hours of greatest usable irradiance for the sintering process.

A control system is then built around the final architecture, splitting the workload across three nodes on a self-contained local WiFi network. An ESP32-CAM acts as both access point and the loop's sensor: it grabs the solar images and serves the centroid error vector over HTTP. A second ESP32, connected as a client, handles estimation and control: it integrates the error into the platform's absolute angle, solves the inverse kinematics by numerical bisection, and drives the actuators in absolute position, sidestepping the drift typical of incremental schemes. A Python console on a PC works as a telemetry station, watching the loop without taking part in it. The tracker therefore runs fully standalone, with neither an internet connection nor any dependency on the PC.

In short, this work shows that the classical Gough-Stewart platform cannot kinematically sustain uninterrupted solar tracking at the studied latitude. It puts forward and theoretically validates an optimised hybrid alternative that pairs closed-loop linkages with visual feedback to deliver the rotational range required for efficient, stiff operation throughout the full solar cycle. Combining this mechanical design with a distributed control loop yields a self-sufficient, robust system that scales naturally to any scenario where pointing accuracy is critical.
\end{abstract}
